\documentclass{scrartcl}
\usepackage{cmap}
\usepackage{fontspec}
\usepackage{lmodern}
\usepackage{microtype}
\usepackage[english]{babel}
\usepackage{tikz}
\usepackage{minted}
\usepackage{xcolor}
\usepackage{ifthen}
\usepackage{hyperref}
\usepackage{fontspec, unicode-math}
\usepackage[normalem]{ulem}

\setmonofont{FiraCode Nerd Font Mono}[Scale=MatchLowercase]

\usetikzlibrary{positioning}
\usetikzlibrary{calc}

\newcommand*{\keyword}[1]{\textbf{\color{green!65!black}{#1}}}
\newcommand*{\Node}{\keyword{node}}
\newcommand*{\Where}{\keyword{where}}
\newcommand*{\Guard}{\keyword{guard}}
\newcommand*{\Assert}{\keyword{assert}}
\newcommand*{\If}{\keyword{if}}
\newcommand*{\Then}{\keyword{then}}
\newcommand*{\Else}{\keyword{else}}
\newcommand*{\Pre}{\keyword{pre}}
\newcommand*{\Fby}{\keyword{fby}}

\newfontfamily\NHLight[
   ItalicFont     = HelveticaNeueLT Std Thin,
   ]{Helvetica Neue LT Std}
\newcommand\textulf[1]{{\NHLight\textit#1}}
\newcommand{\lustrean}[1]{\scalebox{#1}{%
  \begin{tikzpicture}%
    \node (L) {\Huge\(\mathcal{L}\)};%
    \node at (.6,.25) {ustr};%
    \node at (.6,-.2) {\small ${\exists}\hspace*{-.1em}{\forall}$\hspace*{-.1em}\raisebox{-.05ex}{\textulf{N}}};%
  \end{tikzpicture}%
}}
\newminted[lustre]{text}{escapeinside=@@, linenos}
\newmintinline[lustr]{text}{escapeinside=@@}

\setminted{baselinestretch=.8}

\makeatletter
\def\uwave{\bgroup\markoverwith{\lower3.5\p@\hbox{\sixly\textcolor{red!90!black}{\char58}}}\ULon}
\font\sixly=lasy6
\makeatother
\newcommand*{\error}[1]{\uwave{#1}}

\title{Lustrean}
\subtitle{Lustre + Lean = Lustrean}
\author{Jean \textsc{Caspar} \and Adrien \textsc{Mathieu}}
\date{December 2024}

\begin{document}
\maketitle
\tableofcontents

\pagebreak

\section*{Introduction}
Lustrean is an abstract interpreter of a core subset of Lustre in Lean.  It aims at
\begin{itemize}
\item doing a (shallow) embedding of Lustre into Lean;
\item using precise semantics of the embedded program, that is, not to refuse a program just
  because it cannot statically be determined to be meaningful;
\item abstracting over that semantics to semi-decide whether that program will produce runtime
  errors.
\end{itemize}
Runtime errors that Lustrean can prevent are of two kind:
\begin{itemize}
\item the program being meaningless, it produces a non-value;
\item the program fails to verify certain user-specified properties.
\end{itemize}

Moreover, Lustrean is parametric over the abstract domain used.  We provide simple numerical
domains to be used with a nonrelational domain, but this is entirely decoupled from the rest of
the project.

\subsection*{Getting and testing Lustrean}
Lustrean can be cloned from \url{https://github.com/yucto/lustrean}.  The project is divided into
two module trees: the main module, \texttt{Lustrean}; and a module \texttt{Misc} which contains
definitions and lemmas which are of general purpose.

Besides this, a module \texttt{Example} is provided, that shows how to use Lustrean.  It is
recommended that this module is opened with an editor that supports Lean, so as to have immediate
feedback: as the Lustre code is written, it will be checked, compiled and (abstractly)
interpreted on the fly.  For instance, if you type
% this is supposed to work, with a better version of minted...
% [extrakeywords=node where assert guard]
\begin{lustre}
@\Node@ f(x) = o @\Where@
  o = x * x
@\Assert@
  @\error{o ≥ 0}@
\end{lustre}
the last assert will be quickly highlighted in red, indicating that Lustrean has failed to
automatically prove that the said assert will never fail.  If you change the assignment of
\texttt{o} to the following, which is seemingly equivalent but more verbose, you will see that
the error will disappear.
\begin{lustre*}{firstnumber=2}
  o = @\If@ x ≥ 0 @\Then@ x * x @\Else@ x * x
\end{lustre*}
This example will be further explained later on.


\section{Lustre embedding}
Lustrean embeds a core subset of Lustre, enriched with the possibility to declare assertions and
assumptions at a node level.  That is, in addition to local bindings, a node can introduce
assumptions with the \lustr{@\Guard@} keyword (synonyms in other contexts might be
\texttt{assume}, \texttt{require} or pre-conditions, although in Lustrean assumptions can speak
of any variable, not just input variables), and assertions with the \lustr{@\Assert@} keyword
(synonyms are \texttt{ensure}, \texttt{provide} or post-conditions).

\subsection{Syntax}
The subset of Lustre that is accepted by Lustrean is simply typed, with two types: (signed)
integers, and boolean.  However, boolean values can only be used to express guards, assertions
and in used in \texttt{if} statements.  The exact syntax is entirely specified in
\path{Lustrean/Elaboration/Syntax.lean}.

\subsection{Type system}
The type system of Lustrean is extremely simple.  All expressions are (signed) integers, except
for boolean expressions which can only occur in positions where they are syntactically expected.
The type system hasn't been improved with extensions, such as incorporating a clock system, since
the point of doing abstract interpretation is to widen the number of programs for which we can
statically ensure the absence of runtime errors.

All in all, this means that we don't even need a type checking pass, as all the type checking /
inference that we need to do is encoded in the syntax.

\subsection{Reification}
The parser hooks into Lean's own parser, so it actually happens at compile time.  The reification
step turns that syntax tree into a Lean value, that can be manipulated by Lean code.  Since the
meta-language of Lean is Lean itself, this is a very simple step, as the syntax tree is already a
Lean value: it merely does a tree traversal, producing a tree whose type is an ADT encoding
exactly Lutrsean programs.

This is further simplified by the fact that all the analysis happens at compile-time (with
respect to Lean's compilation process), so we don't even need to quote the resulting tree.

\subsection{Inlining}
After having reified a Lustrean program, we know it is syntactically correct (and thus type
correct), but we yet have to check scopes.  We want to avoid having to check at some point that
all the variables are correctly bound, and then transporting that proof through all the
compilation steps until we actually need it to resolve the variables.

For variables bound to nodes, used in node calls, we avoid this problem entirely by inlining them
right after we have resolved them.  This means that either a variable is well-scoped, and it will
be replaced by its definition, or it is not and this will produce an error.

The inlining ensures that every assumption an inlined node makes about its formal parameters is

\subsection{Indicisation}
The remaining variables are of two kind: input variables, or locally bound variables.  Clearly,
we cannot inline input variables, so, instead, we will perform a de Bruijn-like transformation,
replacing every variable with an index, indicating where to find the definition of the variable.

The main differences with a de Bruijn transformation are:
\begin{itemize}
\item that bindings cannot be nested.  Indeed, local bindings are mutually recursive, and need
  not to be statically schedulable;
\item indices are enriched with a proof that we are within bounds.
\end{itemize}
The fact that all bindings are at toplevel means that, after this step has been performed, we end
up with an array of expressions, representing the content of each locally bound variable.  Each
variable is now either an index of an input variable, or an index of a locally bound variable.

\subsection{Normalization}
The normalization step does several code transformations at once.  It is, however, simpler to
explain them individually,
\begin{itemize}
\item Firstly, it transforms \lustr{e@${}_1$@ @\Fby@ e@${}_2$@} into
  \lustr{e@${}_1$@ -> @\Pre@ e@${}_2$@} (the two operators \lustr{->} and \lustr{@\Pre@} will be
  compiled independently);
\item it ensures that, in every expression of the form \lustr{@\Pre@ e}, $e$ is a locally
  bound variable, by creating a new binding \lustr{v = e} and replacing $e$ with $v$ if
  necessary;
\item it introduces a new variable, $n$, which counts the number of times the node has been
  activated.  This is useful to replace \lustr{e@${}_1$@ -> e@${}_2$@} by
  \lustr{@\If@ n = 0 @\Then@ e@${}_1$@ @\Else@ e@${}_2$@};
\item it ensures that every \lustr{@\If@ ... @\Then@ ... @\Else@ ...} expression happens at the 
  top of an expression.  The reason to do so is that the language we will compile Lustrean to
  doesn't have \texttt{if} expressions, only \texttt{if} statements, so we will compile a
  binding \lustr{v = @\If@ c @\Then@ e@${}_1$@ @\Else@ e@${}_2$@} into
  \lustr{@\If@ c @\Then@ v = e@${}_1$@ @\Else@ v = e@${}_2$@};

\item for every locally bound variable \(v\), it creates a variable \(v_\textrm{pre}\), (which in
  ZRun is called ``last'') whose purpose is to store the persistent state.  Then, we replace
  every \lustr{@\Pre@ v} with \lustr{v@${}_\textrm{pre}$@};
\end{itemize}

\subsection{Compilation}
After having normalized the program, we compile it to a Control Flow Graph (CFG onwards) of a
simple imperative programming language (defined in \path{Lustrean/Imp.lean}, which we'll call Imp
onwards), which essentially has two statements: an assignment statement, and a guard statement.
The \lustr{@\Guard@} statement is akin to an \lustr{@\If@} statement, but is formulated as a
  statement that blocks computation if its condition is not met because the generated CFG is not
  deterministic (this will be further explained in the Compilation Scheme section).

This is mostly a virtual step, as, in the code, it is merged with the step that translates a CFG
expressed in the concrete semantics into a CFG in the abstracted semantics.  This is important
because, as will be further explained in the Compilation Scheme section, this allows producing a
much simpler abstracted CFG, by exploiting the abstract interpreter itself to perform some
operations that our concrete program should do explicitly otherwise.

However, this also has a major drawback, which is that at no point the concrete program is
explicitly built.  If we wanted to prove soundness theorems for our abstract interpreter, we
should do it with respect to a concrete semantics.  In fact, a concrete machine does not provide
a fixpoint computation operation, as our abstract interpreter does, and since our compilation
scheme exploits this feature, proving a soundness theorem would additionally require to either
avoid using that primitive fixpoint operation, or also proving a fixpoint lifting theorem that
states that computing a fixpoint in the language, or lifting it in the abstract interpreter ends
up with the same result.

\subsection{Semantics}
The semantics of Lustre that Lustrean provides are akin to those of ZRun (modulo the fact that
ZRun implements more features).  However, they are simplified in the domain used to find a
fixpoint: ZRun has two special, ``undefined'' values, bottom and nil.  The former represents
computation that haven't produced a value yet (and which might never produced a value), while the
latter represents values created by \lustr{@\Pre@}.

In the concrete semantics of Lustrean, instead, these two ``undefined'' values are merged.  The
one value that corresponds to these two values is called, in Lustrean, nil.  This makes the
concrete semantics easy to explain: initialize every (non-input) variable to nil, and propagate
computations.  This applies uniformly for ``last'' variables, and for locally bound variables.

This nil value represents both an undefined value that we are not supposed to use (ie. reading a
\lustr{@\Pre@ e} value at the first step) and divergence/non-causality.  These are two behaviors
that are considered bad, and that Lustrean tries to prove never happen, so there is no real need
to distinguish these two.  However, the concrete semantics cannot actually diverge: computing a
Lustre step always terminates.

These should not be confused with extra ``special'' points added in the abstract domain, $\bot$
and $\top$.  The former represents unreachable expressions (\emph{not} diverging expressions,
since expressions in Lustre cannot diverge anyways), while the latter represents an expression
that might have any value.  While producing a nil (which corresponds, in ZRun, to either bottom
or nil) is considered as bad behavior, producing a $\bot$ or a $\top$ is not bad behavior.

\section{Abstract interpretation}
\subsection{The interpreter}
The abstract interpreter takes, as input, a program (whose form is a CFG backed by simple
imperative statements).  It executes that program, and outputs the state of the abstract machine
at the end of the program.  To do so, it actually computes the state at each vertex (and each
arc) of the CFG, by translating the CFG into an operator over the domain passed as argument, and
computing its fixpoint.

We procede by iteration: we start with every environment to $\bot$, meaning that they are
unreachable, except for the initial node, whose environment is $\top$, meaning that initially any
value might be present in memory. Then we iterate the following steps until the environments in
the node stabilize:
\begin{itemize}
\item we iterate on every arc to assigns to each arc the value of the transfer function of the
  arc applied to the environment of the source node of the arc;
\item we iterate on every node to assign to each node the union of the environments of the arcs
  coming into the node.
\end{itemize}

We couldn't prove that the function iterating until the stabilizition terminates in Lean, as
described in the Domains section, though it may be possible. But Lean allows to annotate as
partial some functions where it can prove that the type of the returned value is inhabited,
ensuring no inconsistency can happen.  In the interpreter, all array accesses are proven correct.

In order to do this, a lot of bookkeeping is involved with dependent types.  For example, we use
two intermediate structures in order to convert the first representation of the CFG, which is
presented as a list of nodes, a nodes consisting in a list of pairs with an instruction and the
index of the exit node of the instruction, and the second representation, which consists of a
list of nodes and arcs, a node containing a list of ``indices into the list of arcs'',
representing the instructions going into this node, and an arc consisting of an instruction and
the source and exit node of this instruction, represented as ``indices into the list of nodes''.

\subsection{Compilation scheme}
The compilation scheme should encode the coiterative semantics within Imp.  Essentially, this
means producing code in that matches correspond to the source code of ZRun, but instantiated with
a particular Lustre program.  This would involve implementing a fixpoint computation in Imp, and
executing over and over again in our outermost fixpoint computation.

However, fixpoint computations can be flattenend, so rather than doing that, we can simply forget
about every condition in the loop computing the innermost fixpoint.  Indeed, \autoref{fig:cfg}
shows the generated CFG, assuming \(\mathcal{L}\) is the set of locally bound variables,
\(\mathcal{I}\) the set of input variables, \(\mathcal{G}\) the set of guards, and
\(\mathcal{A}\) the set of assertions.
% Deduplication of fixpoint computation
\begin{figure}[H]
  \centering
  \begin{tikzpicture}[thick, 1/.style={->}, 2/.style={->, dashed}]
    \node (0) {$\bullet$};

    \foreach[count=\n from 0] \lbl/\style/\right in {
      {$n$ \leftarrow{} $0$}/1/{},%
      {$v$ \leftarrow{} nil}/1/{$\$v\in\mathcal{L}$},%
      {$v_\textrm{pre}$ \leftarrow{} $v$}/1/{$v\in\mathcal{L}$},%
      {$v$ \leftarrow{} $[-\infty, \infty]$}/1/{$v\in\mathcal{I}$},%
      {guard $g$}/1/{$g\in\mathcal{G}$},%
      {$v$ \leftarrow{} nil}/1/{$v\in\mathcal{L}$},%
      {code}/2/{},%
      {assert $a$}/1/{$a\in\mathcal{A}$}%
    } {
      \edef\m{\the\numexpr \n + 1\relax}
      \node[below=of \n] (\m) {$\bullet$};
      \draw[\style] (\n) -- node[right] (h\n) {\lbl} (\m);
      \node[anchor=west] at (h\n -| 3,0) {\right};
    }
    \draw[->] (7) edge[bend left=30] node[left] {skip} (6);
    \draw[->] (7) edge[bend left=90] node[left] {$n$ \leftarrow{} $n + 1$} (2);
  \end{tikzpicture}
  \caption{Generated CFG.  ``code'' refers to executing an assignment for each equation,
    optionally branching if there is an \texttt{if} statement.}\label{fig:cfg}
\end{figure}

Note that this program is non-deterministic: after having ``executed'' the equations, we can
non-deterministically refine once more the current iteration, or skip to the next iteration.

\section{Domains}

The domains used in abstract interpretation are bounded lattice which supports arithmetical
operations, assignments and guards, where a guard is a restriction of the environment to the
values satisfying some predicate. Domains should also come with a widening and a narrowing
operators. A widening operator is an overapproximation of the union of environment, which ensures
we reach a post-fixpoint of the transfer function in finite time. A narrowing is an
overapproximation of the intersection, which allows to improve the precision while staying above
the fixpoint.

\subsection{Non-relational domain}
Here, we only used non relational domains: environments are functions from the set of variables
to a so-called value domain, representing an approximation of a value. In the implemention, we
used vectors rather than functions for more efficiency. Relational domains are far more powerful
than non relatonial domains, but we didn't have the time to implement one. Also, our non
relational domains are reduced, meaning that there is a single representation for all
environments which contain at least one bottom value.

We implemented two value domains: integers with a top and a bottom element, and the intervals. We
also implemented a domain combinator which freely adds a nil value, representing an unknown
value. We use this value to represent non initialized pre expressions. In the new domain, values
are a pair of a value from the old domain et a boolean representing the fact that this value may
be nil. Therefore nil is incomparable with values from the old domain, except with $\bot$. Nil is
greater than $\bot$.

Domains have been fully formalized in Lean, except for the fact that widening and narrowing
stabilizes in finite time: this is would be very painful to prove, since our widening depends on
the list of constants in the program, and may require classical logic, thus being difficult to
extract.

\subsection{Working examples}
The interval non-relational domain is appropriate to prove inequalities between an expression,
and a constant.  Consider, for instance, the following examples for which Lustrean can
successfully disprove runtime failure.
\begin{lustre}
@\Node@ truncate(x) = o @\Where@
  @\Guard@
    x ≥ 0
  @\Where@
    y = x @\Fby@ y + 1
    o = @\If@ y > 3 @\Then@ 3 @\Else@ y
  @\Assert@
    0 ≤ o
    o ≤ 3
\end{lustre}
and
\begin{lustre}
@\Node@ counter() = o @\Where@
  o = 0 @\Fby@ 1 + o
@\Assert@
  o ≥ 0
\end{lustre}
However, it can fail at even the simplest assertion.
\begin{lustre}
@\Node@ embarassing(x) @\Where@
@\Assert@
  x = x
\end{lustre}
If we simplify \(x\) from both sides of this expression, we get the condition \(0 = 0\), which
Lustrean is capable of proving.  Hence, adding a simplifier might greatly increase the precision
of Lustrean.  We will see in the next section an other example where doing a seemingly trivial
transformation (which, this times, apparently complexifies the expression) also helps Lustrean.

\subsection{Non-working examples}
According to our tests, non-relational domains struggle to prove the properties we are interested
in.  However, we will show that it is still possible, in some cases, to recover some kind of
relational-like analysis.

Consider, for instance, the very trivial example
\begin{lustre}
@\Node@ u() = o @\Where@
  o = 0 @\Fby@ o
\end{lustre}
Lustrean fails to automatically prove that this program will not fail at runtime, the reason
being that it will get compiled to
\begin{lustre*}{firstnumber=2}
  o = @\If@ n = 0 @\Then@ 0 @\Else@ o@${}_\textrm{pre}$@
\end{lustre*}
and the non-relational domain cannot express conditions such as ``$n > 0$ implies
$o_\textrm{pre}$ is well-defined''.  Since it is clear that not being able to reason about
whether $n=0$ greatly hinders the power of our analysis, our abstract interpreter would likely be
much more practical if we implemented a relational domain, for instance, by partitioning our
non-relational domain with respect to $n = 0$, that is, partition according to the following
domain:
\begin{center}
  \begin{tikzpicture}
    \node (top) at (1,1) {$\top$};
    \node (eq) at (0,0) {$n = 0$};
    \node (neq) at (2,0) {$n\neq 0$};
    \node (bot) at (1,-1) {$\bot$};

    \draw (bot) -- (eq);
    \draw (bot) -- (neq);
    \draw (eq) -- (top);
    \draw (neq) -- (top);
  \end{tikzpicture}
\end{center}

However, there is actually a way to recover a more precise analysis without much work, by
unrolling the outer loop.  This way, the environment where $n = 0$ and where $n > 0$ end up being
in two different nodes of the graph.  And, indeed, if one actually tries to run the previous
example, Lustrean will successfully not report any error, because it unrolls one iteration of the
loop.  To see the effect of loop unrolling, one can change the value of the constant
\texttt{unroll\_loop}, in \path{Lustrean/Elaboration/Compile.lean:46}.

An other interesting example is the following
\begin{lustre}
@\Node@ f(x) @\Where@
  o = x * x
@\Assert@
  o ≥ 0
\end{lustre}
It is easy to convince oneself that this assertion is always true, yet Lustrean fails to prove
it.  Indeed, in the environment where we know that \(x\in[-\infty, \infty]\), \texttt{x * x}
evaluates to \lustr{[-∞, ∞] * [-∞, ∞]}, which in turn evaluates to \lustr{[-∞, ∞]}.  Yet, not
everything is lost: consider the following program instead
\begin{lustre}
@\Node@ f(x) @\Where@
  o = @\If@ x ≥ 0 @\Then@ x * x @\Else@ x * x
@\Assert@
  o ≥ 0
\end{lustre}
This works because we now evaluate \lustr{x * x} in two different environments, where we can
remember the fact that both operands have the same sign even after substituting \lustr{x}.  We
could apply this transform to every multiplication, but we haven't done so.  Doing so could
result in an exponential blowup in the size of the expression.

A more interesting example to analyze is one that exhibits non-statically schedulable behavior:
\begin{lustre}
@\Node@ f(c, z) = x, y @\Where@
  x = @\If@ c = 0 @\Then@ z @\Else@ y
  y = @\If@ c = 0 @\Then@ x @\Else@ z
\end{lustre}
Lustrean also fails to show this example is well-defined.  We haven't found a way to recover the
information while sticking with a non-relational domain.

\end{document}
